\chapter{Chapter05 — Biased Majority Gaussian Limit}
%%% generated by the blueprint pipeline from FABL.Chapter05.BiasedMajorityGaussianLimit
\section{Overview}

Declarations of the Lean module \texttt{FABL.Chapter05.BiasedMajorityGaussianLimit}.

\section{Declarations}

\begin{definition}[Gaussian upper right quadrant]
\lean{FABL.gaussianUpperRightQuadrant}
\statementsource{Analysis-of-Boolean-Functions-by-Ryan-ODonnell}{
  pdf-ffbe9c7b3306-p120-b003,
  pdf-ffbe9c7b3306-p120-b004,
  pdf-ffbe9c7b3306-p274-b002
}
\label{FABL.gaussianUpperRightQuadrant}
\leanok
\uses{FABL.CorrelationPlane, FABL.correlationFirstCoordinate, FABL.correlationSecondCoordinate}
For $t\in\mathbb R$, the Gaussian upper right quadrant is the set of pairs
$(x,y)\in\mathbb R^2$ such that $t<x$ and $t<y$.
\end{definition}

\begin{theorem}[Openness of the upper-right quadrant]
\lean{FABL.isOpen_gaussianUpperRightQuadrant}
\label{FABL.isOpen_gaussianUpperRightQuadrant}
\leanok
\uses{FABL.continuous_correlationFirstCoordinate, FABL.continuous_correlationSecondCoordinate, FABL.gaussianUpperRightQuadrant}
\difficulty{1}
For each real number $t$, the set of points $(x,y)$ in the Euclidean plane satisfying
$t<x$ and $t<y$ is open.
\end{theorem}

\begin{theorem}[Boundary of an upper right quadrant]
\lean{FABL.frontier_gaussianUpperRightQuadrant_subset}
\label{FABL.frontier_gaussianUpperRightQuadrant_subset}
\leanok
\uses{FABL.CorrelationPlane, FABL.continuous_correlationFirstCoordinate, FABL.continuous_correlationSecondCoordinate, FABL.correlationFirstCoordinate, FABL.correlationSecondCoordinate, FABL.gaussianUpperRightQuadrant}
\difficulty{4}
Let $t\in\mathbb{R}$, and consider the region in the real coordinate plane where both
coordinates are greater than $t$. Every point on its frontier has at least one
coordinate equal to $t$.
\end{theorem}

\begin{theorem}[Nullity of coordinate thresholds]
\lean{FABL.correlatedGaussianMeasure_coordinateThresholds_null}
\label{FABL.correlatedGaussianMeasure_coordinateThresholds_null}
\leanok
\uses{FABL.CorrelationPlane, FABL.continuous_correlationFirstCoordinate, FABL.continuous_correlationSecondCoordinate, FABL.correlatedGaussianMeasure, FABL.correlatedGaussianMeasure_map_firstCoordinate, FABL.correlatedGaussianMeasure_map_secondCoordinate, FABL.correlationFirstCoordinate, FABL.correlationSecondCoordinate}
\difficulty{4}
Let $-1\leq \rho\leq 1$ and $t\in\mathbb R$. For the probability measure on the plane
obtained by taking independent standard Gaussian variables $U,V$ and mapping them to
$(U,\rho U+\sqrt{1-\rho^2}\,V)$, the set where either coordinate equals $t$ has measure
zero.
\end{theorem}

\begin{theorem}[Nullity of the upper right quadrant boundary]
\lean{FABL.correlatedGaussianMeasure_frontier_upperRight_null}
\label{FABL.correlatedGaussianMeasure_frontier_upperRight_null}
\leanok
\uses{FABL.CorrelationPlane, FABL.correlatedGaussianMeasure, FABL.correlatedGaussianMeasure_coordinateThresholds_null, FABL.frontier_gaussianUpperRightQuadrant_subset, FABL.gaussianUpperRightQuadrant}
\difficulty{1}
Let $-1\leq\rho\leq1$ and $t\in\mathbb R$, and let $(X,Y)$ have the probability
distribution obtained from independent standard Gaussian variables $U,V$ by setting
$X=U$ and $Y=\rho U+\sqrt{1-\rho^2}\,V$. The boundary of the set $\{(x,y):t<x\text{ and
}t<y\}$ has probability zero under this distribution.
\end{theorem}

\begin{theorem}[Normalized correlated sign sum measure upper right tendsto]
\lean{FABL.normalizedCorrelatedSignSumMeasure_upperRight_tendsto}
\statementsource{Analysis-of-Boolean-Functions-by-Ryan-ODonnell}{pdf-ffbe9c7b3306-p120-b003}
\statementsource{mossel-neeman-gaussian-noise-stability}{
  pdf-b53c1a5aa27d-p007-b005,
  pdf-b53c1a5aa27d-p035-b001
}
\label{FABL.normalizedCorrelatedSignSumMeasure_upperRight_tendsto}
\leanok
\uses{FABL.correlatedGaussianMeasure, FABL.correlatedGaussianMeasure_frontier_upperRight_null, FABL.gaussianUpperRightQuadrant, FABL.normalizedCorrelatedSignSumMeasure, FABL.normalizedCorrelatedSignSumMeasure_tendsto_correlatedGaussian}
\difficulty{1}
% dataset-pipeline: stub -- prose pending build_informal_statements.py
Statement pending regeneration.
\end{theorem}

\begin{definition}[Gaussian quadrant probability]
\lean{FABL.gaussianQuadrantProbability}
\statementsource{Analysis-of-Boolean-Functions-by-Ryan-ODonnell}{
  pdf-ffbe9c7b3306-p120-b003,
  pdf-ffbe9c7b3306-p120-b004,
  pdf-ffbe9c7b3306-p274-b002
}
\label{FABL.gaussianQuadrantProbability}
\leanok
\uses{FABL.CorrelationPlane, FABL.correlatedGaussianMeasure, FABL.gaussianUpperRightQuadrant, FABL.standardGaussianUpperQuantile}
For a correlation parameter $\rho\in(-1,1)$ and a level $\beta\in(0,1)$, the Gaussian
quadrant probability is the probability, under the centered bivariate standard Gaussian
distribution with correlation $\rho$, that both coordinates exceed the standard Gaussian
upper quantile at level $\beta$.
\end{definition}

\begin{theorem}[Quantile of the Gaussian upper tail]
\lean{FABL.standardGaussianUpperQuantile_upperTailOpen_biasedMajority}
\label{FABL.standardGaussianUpperQuantile_upperTailOpen_biasedMajority}
\leanok
\uses{FABL.standardGaussianUpperQuantile, FABL.standardGaussianUpperQuantile_upperTail, FABL.standardGaussianUpperTail, FABL.standardGaussianUpperTailOpen, FABL.standardGaussianUpperTail_lt_one, FABL.standardGaussianUpperTail_pos}
\difficulty{3}
For every real number $t$, the standard Gaussian upper quantile at the standard Gaussian
upper-tail probability of $t$ is $t$.
\end{theorem}

\begin{theorem}[First coordinate as a normalized sign sum]
\lean{FABL.correlationFirstCoordinate_normalizedCorrelatedPairSum_eq_normalizedRademacherSum}
\label{FABL.correlationFirstCoordinate_normalizedCorrelatedPairSum_eq_normalizedRademacherSum}
\leanok
\uses{FABL.correlationFirstCoordinate, FABL.correlationFirstCoordinate_normalizedCorrelatedPairSum, FABL.linearForm, FABL.normalizedCorrelatedPairSum, FABL.normalizedRademacherSum}
\difficulty{3}
For every natural number $n$ and every sign vector $x$ of length $n$, the first
coordinate of the normalized correlated sum formed from a pair whose first component is
$x$ equals the normalized Rademacher sum of $x$, namely $(\sqrt n)^{-1}\sum_{i=1}^n
x_i$, with signs interpreted as real numbers.
\end{theorem}

\begin{theorem}[Second coordinate as a normalized sign sum]
\lean{FABL.correlationSecondCoordinate_normalizedCorrelatedPairSum_eq_normalizedRademacherSum}
\label{FABL.correlationSecondCoordinate_normalizedCorrelatedPairSum_eq_normalizedRademacherSum}
\leanok
\uses{FABL.correlationSecondCoordinate, FABL.correlationSecondCoordinate_normalizedCorrelatedPairSum, FABL.linearForm, FABL.normalizedCorrelatedPairSum, FABL.normalizedRademacherSum}
\difficulty{3}
For every natural number $n$ and every pair $x,y\in\{ -1,1\}^n$ of sign vectors, the
second coordinate of their normalized correlated sum equals the normalized Rademacher
sum of $y$, namely $(\sqrt n)^{-1}\sum_{i=1}^n y_i$.
\end{theorem}

\begin{theorem}[Noise stability as a quadrant probability]
\lean{FABL.coe_normalizedCorrelatedSignSumMeasure_upperRight_eq_noiseStability}
\label{FABL.coe_normalizedCorrelatedSignSumMeasure_upperRight_eq_noiseStability}
\leanok
\uses{FABL.CorrelationPlane, FABL.correlatedPairPMF, FABL.correlationFirstCoordinate_normalizedCorrelatedPairSum_eq_normalizedRademacherSum, FABL.correlationSecondCoordinate_normalizedCorrelatedPairSum_eq_normalizedRademacherSum, FABL.gaussianUpperRightQuadrant, FABL.hammingUpperTailIndicator, FABL.isOpen_gaussianUpperRightQuadrant, FABL.noiseStability, FABL.normalizedCorrelatedPairSum, FABL.normalizedCorrelatedSignSumMeasure, FABL.normalizedCorrelatedSignSumMeasure_eq_map_correlatedPairPMF, FABL.normalizedRademacherSum, FABL.pmfExpectation, FABL.pmfExpectation_eq_integral}
\difficulty{5}
Let $\rho\in[-1,1]$, let $t\in\mathbb R$, and let $n$ be a nonnegative integer. The
probability that both coordinates of the normalized sum of $n$ independent correlated
sign pairs exceed $t$ equals the noise stability at correlation $\rho$ of the indicator
that a normalized sum of $n$ independent uniform signs exceeds $t$.
\end{theorem}

\begin{theorem}[Tendsto noise stability hamming upper tail indicator]
\lean{FABL.tendsto_noiseStability_hammingUpperTailIndicator}
\statementsource{Analysis-of-Boolean-Functions-by-Ryan-ODonnell}{
  pdf-ffbe9c7b3306-p120-b002,
  pdf-ffbe9c7b3306-p120-b003,
  pdf-ffbe9c7b3306-p275-b002
}
\label{FABL.tendsto_noiseStability_hammingUpperTailIndicator}
\leanok
\uses{FABL.coe_normalizedCorrelatedSignSumMeasure_upperRight_eq_noiseStability, FABL.correlatedGaussianMeasure, FABL.gaussianQuadrantProbability, FABL.gaussianUpperRightQuadrant, FABL.hammingUpperTailIndicator, FABL.noiseStability, FABL.normalizedCorrelatedSignSumMeasure, FABL.normalizedCorrelatedSignSumMeasure_upperRight_tendsto, FABL.standardGaussianUpperQuantile_upperTailOpen_biasedMajority, FABL.standardGaussianUpperTailOpen}
\difficulty{4}
% dataset-pipeline: stub -- prose pending build_informal_statements.py
Statement pending regeneration.
\end{theorem}

\begin{theorem}[Symmetry under simultaneous negation]
\lean{FABL.correlatedGaussianMeasure_map_neg_biasedMajority}
\label{FABL.correlatedGaussianMeasure_map_neg_biasedMajority}
\leanok
\uses{FABL.CorrelationPlane, FABL.charFun_correlatedGaussianMeasure, FABL.correlatedGaussianMeasure, FABL.correlationQuadraticForm}
\difficulty{4}
Let $-1\leq\rho\leq1$, and let $(X,Y)$ have the probability distribution obtained from
independent standard Gaussian variables $U,V$ by setting $X=U$ and $Y=\rho
U+\sqrt{1-\rho^2}\,V$. The distribution of $(-X,-Y)$ equals the distribution of $(X,Y)$.
\end{theorem}

\begin{theorem}[Symmetry of Gaussian quadrant probabilities]
\lean{FABL.correlatedGaussianMeasure_upperRight_neg_eq_lowerLeft}
\label{FABL.correlatedGaussianMeasure_upperRight_neg_eq_lowerLeft}
\leanok
\uses{FABL.CorrelationPlane, FABL.correlatedGaussianMeasure, FABL.correlatedGaussianMeasure_coordinateThresholds_null, FABL.correlatedGaussianMeasure_map_neg_biasedMajority, FABL.correlationFirstCoordinate, FABL.correlationSecondCoordinate, FABL.gaussianUpperRightQuadrant, FABL.isOpen_gaussianUpperRightQuadrant}
\difficulty{5}
Let $\rho\in[-1,1]$, and let $(X,Y)$ have the centered bivariate Gaussian law with unit
variances and correlation $\rho$. For every $t\in\mathbb R$, the probability that $X>-t$
and $Y>-t$ equals the probability that $X\le t$ and $Y\le t$.
\end{theorem}

\begin{theorem}[Gaussian lower-left quadrant probability]
\lean{FABL.gaussianQuadrantProbability_one_sub_upperTail}
\label{FABL.gaussianQuadrantProbability_one_sub_upperTail}
\leanok
\uses{FABL.CorrelationPlane, FABL.correlatedGaussianMeasure, FABL.correlatedGaussianMeasure_upperRight_neg_eq_lowerLeft, FABL.correlationFirstCoordinate, FABL.correlationSecondCoordinate, FABL.gaussianQuadrantProbability, FABL.standardGaussianUpperQuantile_one_sub, FABL.standardGaussianUpperQuantile_upperTailOpen_biasedMajority, FABL.standardGaussianUpperTailOpen}
\difficulty{3}
Let $-1<\rho<1$ and $t\in\mathbb R$. If $\beta=1-\mathbb P(Z>t)$ for a standard normal
random variable $Z$, then the probability that both coordinates of a centered bivariate
standard normal pair with correlation $\rho$ are at most $t$ equals the probability that
both coordinates are greater than the upper quantile corresponding to $\beta$.
\end{theorem}
